\section{Módulos adicionales y componentes}

Este bloque documenta todo lo demás que existe al \textbf{27-11-2026}. El
software no está terminado: la próxima revisión de documentación actualizará
este y los demás bloques con lo que se implemente después.

\subsection{Lo agregado entre el 18 y el 27 de noviembre}

\begin{itemize}[leftmargin=*,itemsep=0.2em]
\item \textbf{Ficha de salud del personal} en \codigo{/FichasMedicas}
  (tablas y auditoría de los pasos 25--26), reservada a UTP y Dirección.
\item \textbf{Observaciones de situación familiar} (paso 21), separadas del
  área de salud.
\item \textbf{Notas de retiro} (pasos 17 y 24): las escriben Secretaría, UTP
  y Convivencia.
\item \textbf{Cursos sin letra}: parvularia único (paso 18) y luego
  pre-kínder y kínder separados (paso 23), con la decisión del colegio.
\item \textbf{Auditoría legible} para Dirección e Inspectoría (paso 19) y
  \textbf{firmante persistido} en el certificado (paso 20).
\item \textbf{Conservación}: solo el resumen técnico de alertas se poda a 24
  meses (paso 22); expediente y auditoría no se borran.
\item \textbf{ADMIN\_TI sin acceso a datos de alumnos} (paso 27) y
  \textbf{\codigo{EMITIR}} admitido en \codigo{registros\_acceso} (paso 28):
  antes la emisión de certificados no dejaba rastro.
\item \textbf{Medicación habilitada} (\codigo{Medicacion:Habilitada}) por
  orden expresa del 25-11-2026, con el protocolo escrito aún pendiente.
\end{itemize}

\subsection{Módulo Auth: identidad neutral (\codigo{mpaz\_auth})}

Esquema neutral de identidad del colegio: no pertenece a ningún sistema, es
el contrato que todos leen y ninguno edita directamente. Definido en
\codigo{database/modulos/auth/pasos/10-ddl-auth.sql} (paso 10, convención de
pasos numerados desde 10).

\begin{longtable}{@{}p{0.26\textwidth}p{0.66\textwidth}@{}}
\toprule
\textbf{Tabla} & \textbf{Contenido}\\
\midrule
\codigo{persona} & Una fila por persona: \codigo{id} UUID (= el
  \codigo{sub} de los tokens SSO, se conserva aunque cambie el correo),
  \codigo{correo} único (índice case-insensitive), \codigo{nombre},
  \codigo{activo}. Sembrada desde \codigo{funcionario} con
  \codigo{datos-prueba/10-siembra-persona.sql} (solo desarrollo).\\
\codigo{credencial} & Reservada para fase 2 (mudar aquí los hashes de
  Identity). Hoy vacía a propósito: las claves siguen en
  \codigo{mpaz\_matricula.AspNetUsers}.\\
\codigo{OpenIddict*} & Aplicaciones, autorizaciones, scopes y tokens del
  emisor. Las crea la migración EF del \codigo{SsoDbContext}, no los
  scripts SQL.\\
\bottomrule
\end{longtable}

Reglas del módulo: las altas, bajas y cambios de correo los hace
Administración TI; las aplicaciones solo leen \codigo{persona} (por
\codigo{id} o \codigo{correo}); ningún sistema guarda claves fuera de aquí
(fase 2) ni duplica nombre/correo como verdad propia;
\codigo{funcionario.sso\_sujeto = persona.id}.

En código: \codigo{SsoDbContext} (\codigo{Infrastructure/Persistence}, esquema
por defecto \codigo{mpaz\_auth}, \codigo{UseOpenIddict()}) y
\codigo{EmisorSso} (\codigo{Web/Seguridad}): endpoints
\codigo{/connect/authorize} y \codigo{/connect/token}, scopes
\codigo{openid}, \codigo{profile}, \codigo{email} y \codigo{mpaz\_api},
flujos authorization-code y refresh-token. Solo desarrollo: usa certificado
de desarrollo y \codigo{DisableTransportSecurityRequirement} en http local;
en producción el arranque se detiene hasta configurar el certificado real
de firma. \codigo{SembrarAsync} registra/actualiza el cliente de prueba
desde user-secrets (\codigo{Sso:ClientePruebaId/Secreto/Retorno}) en cada
arranque de desarrollo, sin recrear la base. Prototipo de prueba:
\codigo{scripts/probar-sso.py} (requiere \codigo{MPAZ\_DB\_CONEXION}).

\subsection{Módulo Asistencia: semilla DDL (\codigo{mpaz\_asistencia})}

Esquema traído de la base \codigo{mpaz\_asistencia} de Neon del equipo de
asistencia, consolidado como semilla en
\codigo{database/modulos/asistencia/pasos/10-ddl-asistencia.sql}: 11 tablas
y 9 enums. \textbf{No tiene flujo en el portal de matrícula} (MAT-15
pendiente, sin datos de prueba aún): se documenta solo su esquema.

Tablas: \codigo{alumnos} (RUN único, curso, estado activo/inactivo/retirado),
\codigo{cursos}, \codigo{asignaturas}, \codigo{asignaciones\_docentes}
(profesor, curso, asignatura), \codigo{sesiones\_clase} (borrador/confirmada/
anulada, modalidad por clase/diaria), \codigo{asistencias\_clase} (presente,
atrasado, justificado, no presente; origen manual/RFID),
\codigo{asistencias} (jornada RFID: una fila por alumno y día, \codigo{fecha}
generada en America/Santiago, entrada/salida por totem),
\codigo{tarjetas} (UID RFID, estados disponible/asignada/bloqueada/
extraviada), \codigo{tarjeta\_historial} (bitácora, \codigo{alumno\_id}
nunca se nulifica), \codigo{totems} y \codigo{usuarios} (administrador/
profesor).

Reglas relevantes: unicidad alumno/día en \codigo{asistencias} y
sesión/alumno en \codigo{asistencias\_clase}; trigger
\codigo{ck\_alumno\_en\_curso} valida que el alumno pertenezca al curso de
la sesión; restricción de exclusión impide sesiones solapadas por curso
(no anuladas); \codigo{asistencias} y \codigo{asistencias\_clase} son
independientes a propósito y el único puente es \codigo{alumnos}.

\subsection{Web service externo de registro}

\codigo{RegistroWebService} (\codigo{Web/Servicios}) sincroniza funcionarios
con el web service externo: \codigo{POST \{WebService:BaseUrl\}/api/registro/}
con cabecera \codigo{Authorization: Token ...} y cuerpo JSON
\codigo{\{nombre, correo, contrasena\}}. La base por defecto es
\codigo{https://api.empaz.cl/api/registro/}; el token vive en
\codigo{WebService:Token} (user-secrets o entorno, nunca en el repo).
\codigo{Configurado} indica si hay base y token. Si el web service está caído,
no configurado o rechaza, el método devuelve \codigo{false} y el funcionario
ya quedó cargado en Neon: la base del portal es el respaldo.

\subsection{Estructura por módulos y componentes de despliegue}

La base de datos se organiza en \codigo{database/}: \codigo{comun/pasos/}
(00 rol de aplicación, 02 Identity, compartidos), \codigo{modulos/matricula/}
(implementado), \codigo{modulos/asistencia/} y \codigo{modulos/auth/}
(semillas, paso 10), \codigo{modulos/\_plantilla-modulo/} (esqueleto para
futuros módulos como notas o apoyo, no es un módulo real) y el ejecutor
\codigo{aplicar.sh} (\codigo{--modulo} repetible, \codigo{--todos},
\codigo{--neon}). Cada módulo usa \codigo{instalacion/} (bases nuevas),
\codigo{pasos/} (mantenimiento) y \codigo{datos-prueba/} (solo desarrollo,
excluidos del release).

Componentes de despliegue: \codigo{mpaz-matricula/deploy/}
(\codigo{build-release.sh}, \codigo{mpaz-matriculas.service},
\codigo{mpaz-matriculas.env.example}) y \codigo{scripts/publicar-paquete.sh}
(release \codigo{linux-x64} con SHA-256). Las carpetas
\codigo{BackgroundServices/} y \codigo{Endpoints/} de Web existen pero están
vacías: el único proceso en segundo plano, la revisión de alertas cada 15
minutos, vive en \codigo{Servicios/AlertasProgramadas.cs} y se registra como
\codigo{BackgroundService} en \codigo{Program.cs}.

\nota{Corte de esta documentación: 27-11-2026. Todo lo descrito aquí es lo
que existe en el repositorio a esa fecha; lo que se implemente después se
incorporará en la próxima revisión.}
